\section{Dataset and Preprocessing}

\subsection{Dataset Description}
This study focuses on seven Indonesian banking stocks: BBCA, BBNI, BBRI, BBTN, BDMN, BMRI, and BNGA. Daily market data include price and valuation-related variables, while macroeconomic variables include inflation-related indicators, policy rate, exchange rate, and market index features.

The target variable is the daily closing price. The analysis period follows a chronological split to avoid data leakage.

\subsection{Feature Selection}
The baseline feature set (\texttt{Without\_Macro}) includes market micro and technical variables available in the bank-level dataset. The macro-augmented regime (\texttt{With\_Macro}) extends this set with external variables:
\begin{itemize}
    \item inflation rate and CPI-derived indicators,
    \item BI policy rate,
    \item USD/IDR exchange rate,
    \item IHSG index level and return proxy.
\end{itemize}

The final feature schema is fixed per regime before model training.

\subsection{Feature Schema for SHAP Analysis}
To support consistent SHAP analysis across model families, feature names are kept identical across banks and seeds, and the same temporal windowing process is used in every run. The \texttt{Without\_Macro} regime uses 10 market/valuation features, while \texttt{With\_Macro} adds 6 external variables (total 16 features). This fixed schema allows direct aggregation of feature attribution statistics for original recurrent and Transformer outputs, and it keeps the recurrent-variant outputs compatible with future SHAP analysis.

\subsection{Data Preprocessing}
We apply chronological sorting, forward-fill for missing observations, and standard scaling for both features and target during model training. Sequence windows are then constructed with a fixed lookback length.

Time split protocol:
\begin{itemize}
    \item Training: 2014-01-02 to 2019-12-31
    \item Validation: 2020-01-02 to 2021-12-31
    \item Testing: 2022-01-02 to 2024-12-30
\end{itemize}

This split reflects strict temporal validation and supports fair out-of-sample evaluation.
